\textbf{Commun à tous les candidats}

\medskip














\textbf{Question 1 :}

On considère la fonction $g$ définie sur $]0~;~+\infty[$  par $g(x)= x^2+ 2x - \dfrac{3}{x}$.







Si $T$ est la tangente à la courbe représentative de $g$ au point d'abscisse 1, on sait que :

$M(x~;~y) \in T \iff y - g(1) = g'(1)(x - 1)$.

\starredbullet $g$ est dérivable sur $]0~;~+\infty[$ et sur cet intervalle : $g'(x) = 2x + 2 + \dfrac{3}{x^2}$

On a donc $g'(1) = 2 + 2 + \dfrac{3}{1} = 7$ ;

\starredbullet D'autre part $g(1) = 1^2 + 2 \times 1 - \dfrac{3}{1} = 0$.

On a donc : $M(x~;~y) \in T \iff y - 0 = 7(x - 1) \iff y = 7x - 7$. Réponse \textbf{a.}

\medskip

\textbf{Question 2 :}









Pour $n$ assez grand, on a $n \ne 0$, donc $v_n = \dfrac{3}{1 + \frac{2}{n}}$

Comme $\displaystyle\lim_{n \to + \infty} \dfrac{3}{n} = 0$, on a $\displaystyle\lim_{n \to + \infty} = 3$. Réponse \textbf{b.}

\medskip

\textbf{Question 3 :}








Si $X$ est la variable aléatoire correspondant au nombre de boules noires tirées, l'expérience correspond à une loi de Bernoulli  et $X$ suit une loi binomiale de paramètres $n = 10$ et $p = 0,6$.

On sait que $p(X = 4) = \binom{10}{4} \times 0,6^4 \times \left(1 - 0,4^6\right) \approx \np{0,11147} \approx \np{0,1115}$ à $10^{-4}$ près. Réponse \textbf{c.~~}
\medskip

\textbf{Question 4 :}









On a pour $x \ne 0$, \, $f(x) = x\left(3\dfrac{\e^x}{x}  - 1\right)$.

Or on sait que $\displaystyle\lim_{x \to + \infty} \dfrac{\text{e}^x}{x} = + \infty$, donc 
$\displaystyle\lim_{x \to + \infty} \dfrac{\text{e}^x}{x} - 1 = + \infty$ et par produit de limites
 
$\displaystyle\lim_{x \to + \infty} f(x) = + \infty$ Réponse \textbf{b.}
\medskip

\textbf{Question 5 :}














Le nombre de codes différents est donc $36^8 = \np{2821109907456}$.

Il faut donc $\dfrac{\np{2821109907456}}{10^8} \approx \np{28211}$~(s) ou $\approx 478,2$~(min) ou $\approx 7,8$~(h). Réponse \textbf{b.}

\bigskip
